% =====================================================================
%  Bölüm 1 — Mantık Nedir, İspat Asistanı Nedir?
% =====================================================================
\chapter{Mantık Nedir, İspat Asistanı Nedir?}\label{bol:giris}

\begin{bolumhedefleri}
\begin{itemize}
  \item Bir çıkarımın \emph{geçerli} olmasının ne demek olduğunu ve
        geçerlilik ile doğruluk arasındaki farkı,
  \item Mantığın Aristoteles'ten yapay zekâya uzanan kısa tarihini,
  \item Bir ispat asistanının nasıl çalıştığını ve neden güvenilir olduğunu,
  \item \Lean{}'i kurmayı, \lc{\textcolor{leancmd}{\bfseries\#check}} ve
        \lc{\textcolor{leancmd}{\bfseries\#eval}} komutlarını kullanmayı,
  \item İlk teoreminizi yazmayı ve \Lean{}'in hata mesajlarını okumayı.
\end{itemize}
\end{bolumhedefleri}

% ---------------------------------------------------------------------
\section{Mantık nedir?}\label{sec:mantik-nedir}

Mantık, \emph{doğru akıl yürütmenin} bilimidir. Daha kesin söylersek, mantık
hangi sonuçların hangi öncüllerden \emph{zorunlu olarak} çıktığını inceler.
Bu tanımdaki kilit kelime ``zorunlu olarak''dır. Mantık, öncüllerin doğru
olup olmadığıyla ilgilenmez; öncüller doğru \emph{olsaydı} sonucun da doğru
olmak zorunda olup olmadığıyla ilgilenir.

\begin{tanim}{Çıkarım ve geçerlilik}{cikarim}
Bir \textbf{çıkarım}, bir öncüller kümesi ile bir sonuçtan oluşur.
Öncüllerin hepsinin doğru olduğu her durumda sonuç da doğruysa,
çıkarım \textbf{geçerlidir}. Geçerli olup öncülleri de gerçekten doğru
olan bir çıkarıma \textbf{sağlam} çıkarım denir.
\end{tanim}

\begin{ornek}{Geçerlilik ve doğruluk birbirinden bağımsızdır}{gecerlilik}
\begin{center}
\renewcommand{\arraystretch}{1.3}
\begin{tabular}{p{0.44\linewidth} p{0.44\linewidth}}
\textbf{(a) Geçerli ve sağlam}
& \textbf{(b) Geçerli ama sağlam değil} \\
Bütün insanlar ölümlüdür. & Bütün kuşlar uçar. \\
Sokrates bir insandır. & Penguen bir kuştur. \\
\cmidrule(r){1-1}\cmidrule(l){2-2}
O hâlde Sokrates ölümlüdür. & O hâlde penguen uçar. \\
\end{tabular}
\end{center}
İki çıkarımın \emph{biçimi} aynıdır: ``Bütün $A$'lar $B$'dir; $x$ bir
$A$'dır; o hâlde $x$ bir $B$'dir.'' Bu biçim geçerlidir. (b)'deki sonucun
yanlış olması mantığın değil, ilk öncülün suçudur.
\end{ornek}

Bu örnek, modern mantığın temel fikrini gösteriyor: \textbf{geçerlilik,
içerikle değil biçimle ilgilidir}. Sokrates'in ya da penguenin ne olduğunu
bilmemize gerek yoktur; yalnızca çıkarımın iskeletine bakarız. Bu iskeleti
sembollerle yazdığımızda mantık \emph{formal} (biçimsel) mantık olur:
\[
  \frac{\forall x\,(A(x) \to B(x)) \qquad A(s)}{B(s)}
\]
Kesrin üstündekiler öncüller, altındaki sonuçtur. Kitap boyunca bu gösterimi
çok kullanacağız.

\begin{onemlifikir}
Biçimsel mantıkta bir çıkarımın geçerliliğini denetlemek \emph{mekanik} bir
iştir: Sembollerin anlamını bilmeden, yalnızca kurallara bakarak yapılabilir.
Mekanik olan her iş, bir bilgisayara yaptırılabilir. İspat asistanlarının
varlık nedeni budur.
\end{onemlifikir}

% ---------------------------------------------------------------------
\section{Kısa bir tarihçe}\label{sec:tarihce}

\Cref{fig:zaman-cizelgesi}, bu kitapta karşılaşacağımız fikirlerin
tarihsel sırasını gösteriyor. Dikkat edilirse çizelge üç döneme ayrılıyor:
mantığın \emph{doğuşu}, mantığın \emph{matematikleşmesi} ve mantığın
\emph{makineleşmesi}.

\begin{figure}[htbp]
\centering
\begin{tikzpicture}[
    olay/.style={font=\baslikfont\scriptsize, align=center, text width=2.3cm},
    yil/.style={font=\baslikfont\scriptsize\bfseries},
    nokta/.style={circle, fill=#1, inner sep=2.2pt},
  ]
  % Dönem şeritleri
  \fill[ornekbg]    (0,-0.25) rectangle (2.2,0.25);
  \fill[tanimbg]    (2.2,-0.25) rectangle (8.2,0.25);
  \fill[fikirbg]    (8.2,-0.25) rectangle (14.6,0.25);
  \draw[-{Stealth[length=3mm]}, very thick, koyu] (0,0) -- (14.9,0);
  \node[font=\baslikfont\scriptsize\bfseries, text=ornekframe] at (1.1,-3.0) {Doğuş};
  \node[font=\baslikfont\scriptsize\bfseries, text=anarenk]    at (5.2,-3.0) {Matematikleşme};
  \node[font=\baslikfont\scriptsize\bfseries, text=ikincil]    at (11.4,-3.0) {Makineleşme};

  % Olaylar: (x, yıl, metin, yukarı/aşağı, renk)
  \foreach \x/\y/\t/\d/\c in {
      0.6/MÖ 350/Aristoteles:\\kıyaslar/1/ornekframe,
      2.6/1854/Boole:\\düşünce cebiri/-1/anarenk,
      3.6/1879/Frege:\\Begriffsschrift/1/anarenk,
      4.8/1910/Russell–Whitehead:\\Principia/-1/anarenk,
      5.9/1931/Gödel:\\eksiklik/1/anarenk,
      7.0/1935/Gentzen:\\doğal çıkarım/-1/anarenk,
      8.0/1940/Church:\\tip kuramı/1/anarenk,
      9.1/1967/de Bruijn:\\Automath/-1/ikincil,
      10.0/1969/Howard:\\formüller = tipler/1/ikincil,
      11.0/1988/Coquand–Huet:\\CoC/-1/ikincil,
      12.0/2015/Lean/1/ikincil,
      13.0/2021/Lean 4/-1/ikincil,
      14.2/2024+/YZ destekli\\ispatçılar/1/vurgu}
  {
    \node[nokta=\c] at (\x,0) {};
    \draw[\c, thin] (\x,0) -- (\x,{0.85*\d});
    \node[yil, text=\c] at (\x,{1.05*\d}) {\y};
    \node[olay, text width=1.5cm] at (\x,{1.62*\d+0.38*min(\d,0)}) {\t};
  }
\end{tikzpicture}
\caption{Bu kitaptaki fikirlerin kısa tarihi. Yıllar, ilgili eserin
yayımlandığı ya da sistemin duyurulduğu yıllardır; Howard'ın notları 1969'da
dolaşıma girmiş, 1980'de yayımlanmıştır.}
\label{fig:zaman-cizelgesi}
\end{figure}

\begin{tarihce}[Tarihçe: Frege ve kavram yazısı]
Gottlob Frege 1879'da yayımladığı \emph{Begriffsschrift} (Kavram Yazısı)
adlı küçük kitapta, niceleyicileri (\emph{her}, \emph{bazı}) içeren ilk
tam biçimsel mantık dilini kurdu \citep{frege1879}. Frege'nin amacı,
aritmetiğin tamamen mantıktan türetilebileceğini göstermekti. Bu program
Russell paradoksuyla sarsıldı; ancak Frege'nin kurduğu dil, bugün
kullandığımız yüklemler mantığının doğrudan atasıdır.
\end{tarihce}

% ---------------------------------------------------------------------
\section{Neden bilgisayarla ispat?}\label{sec:neden}

Matematikçiler binlerce yıldır ispatlarını kâğıt üzerinde yazıp birbirlerine
denetletiyorlar. Bu sistem şaşırtıcı derecede iyi çalışır, ama kusursuz
değildir. İspatlar uzadıkça ve karmaşıklaştıkça, hakemlerin her adımı
denetlemesi olanaksızlaşır. Yirminci yüzyılın sonunda bu sorun somut olarak
ortaya çıktı:

\begin{itemize}
  \item \textbf{Dört renk teoremi} (1976), ilk kez büyük ölçüde bilgisayar
  hesabına dayanan bir ispatla kanıtlandı. İspatın tamamı daha sonra Coq
  ispat asistanında baştan sona formalize edildi.
  \item \textbf{Kepler sanısı}nın ispatı o kadar karmaşıktı ki hakemler
  doğruluğundan ``yüzde doksan dokuz emin'' olabildiklerini bildirdi. Bunun
  üzerine başlatılan \emph{Flyspeck} projesi, ispatı HOL Light ve Isabelle
  ispat asistanlarında tamamen doğruladı.
\end{itemize}

Bugün ispat asistanları üç farklı amaca hizmet ediyor:

\begin{center}
\begin{tikzpicture}[
    kutu/.style={draw=#1, fill=#1!8, thick, rounded corners=4pt,
                 text width=4.1cm, minimum height=2.4cm, align=center,
                 font=\small},
  ]
  \node[kutu=anarenk] (a) at (0,0)
    {\baslikfont\bfseries\color{anarenk} Matematiği doğrulamak\\[3pt]
     \normalfont Uzun ve karmaşık ispatların her adımını denetlemek};
  \node[kutu=ikincil] (b) at (4.7,0)
    {\baslikfont\bfseries\color{ikincil} Yazılımı doğrulamak\\[3pt]
     \normalfont Derleyicilerin, işletim sistemlerinin, şifreleme
     kodlarının doğruluğunu ispatlamak};
  \node[kutu=vurgu] (c) at (9.4,0)
    {\baslikfont\bfseries\color{vurgu} Öğrenmek ve öğretmek\\[3pt]
     \normalfont Her adımı anında geri bildirimle öğrenmek};
\end{tikzpicture}
\end{center}

Bu kitap üçüncü amaçla başlıyor, birinciye doğru ilerliyor. İkinci amaç
(yazılım doğrulama) kitabın kapsamı dışında kalsa da, öğreneceğiniz her
teknik orada da doğrudan kullanılabilir.

% ---------------------------------------------------------------------
\section{Bir ispat asistanı nasıl çalışır?}\label{sec:nasil-calisir}

Bir ispat asistanına ``bu teorem doğrudur'' dediğinizde, size
\emph{güvenmez}. Sizden bir ispat ister ve o ispatı adım adım denetler.
\Lean{}'de bu süreç iki katmanda gerçekleşir (\cref{fig:lean-katmanlar}).

\begin{figure}[htbp]
\centering
\begin{tikzpicture}[
    katman/.style={draw=#1, fill=#1!10, thick, rounded corners=4pt,
                   minimum width=3.0cm, minimum height=1.5cm, align=center,
                   font=\baslikfont\small},
    ok/.style={-{Stealth[length=3mm]}, very thick, koyu!70},
    aciklama/.style={font=\scriptsize, text=black!65, align=center,
                     text width=3.0cm},
  ]
  \node[katman=ornekframe] (k) at (0,0) {\bfseries Siz\\ \small taktikler, terimler};
  \node[katman=anarenk]    (e) at (4.9,0) {\bfseries Elaboratör\\ \small büyük, akıllı};
  \node[katman=teoremframe](c) at (9.8,0) {\bfseries Çekirdek\\ \small küçük, güvenilir};
  \node[circle, draw=bittiframe, fill=bittibg, very thick, minimum size=1.1cm,
        font=\sembolfont\Large\bfseries\color{bittiframe}] (t) at (12.6,0) {✓};

  \draw[ok] (k) -- node[above, font=\scriptsize]{kaynak kod} (e);
  \draw[ok] (e) -- node[above, font=\scriptsize]{ispat terimi} (c);
  \draw[ok] (c) -- (t);

  \node[aciklama, below=0.25cm of k] {Okunabilir, kısa ispatlar yazarsınız.};
  \node[aciklama, below=0.25cm of e] {Taktikleri çalıştırır, eksikleri
        tamamlar, ayrıntılı bir ispat terimi üretir.};
  \node[aciklama, below=0.25cm of c] {Yalnızca terimin türünü denetler.
        Hata yaparsa bunu \emph{yakalar}.};

  \draw[decorate, decoration={brace, amplitude=6pt, mirror}, thick, black!40]
    ([yshift=-1.75cm]e.south west) -- ([yshift=-1.75cm]e.south east)
    node[midway, below=8pt, font=\scriptsize, text=black!60, align=center]
    {hata yapabilir\\(ama zarar veremez)};
\end{tikzpicture}
\caption{\Lean{}'in iki katmanlı mimarisi. Güven yalnızca küçük çekirdeğe
dayanır.}
\label{fig:lean-katmanlar}
\end{figure}

Bu tasarımın özü şudur: Elaboratör çok büyük ve karmaşık bir programdır,
dolayısıyla içinde hata olabilir. Ama elaboratörün ürettiği her ispat
terimi, ayrıca küçük bir \emph{çekirdek} tarafından denetlenir. Çekirdek
yanlış bir ispatı kabul etmedikçe, elaboratördeki bir hata yanlış bir
teoremin ``ispatlanmasına'' yol açamaz.

\begin{tanim}{de Bruijn ölçütü}{debruijn}
Bir ispat asistanının ürettiği ispatlar, \textbf{küçük ve bağımsız bir
denetleyici} tarafından yeniden doğrulanabiliyorsa, o sistem \textbf{de
Bruijn ölçütünü} sağlar. Böylece sisteme duyulan güven, yalnızca bu küçük
denetleyiciye duyulan güvene indirgenir.
\end{tanim}

\begin{tarihce}[Tarihçe: Automath]
Hollandalı matematikçi Nicolaas Govert de Bruijn, 1967'de başlattığı
\emph{Automath} projesiyle matematiksel metinleri bilgisayarda denetlemenin
ilk ciddi örneğini verdi \citep{debruijn1970}. Automath ekibi 1970'lerde
Landau'nun \emph{Grundlagen der Analysis} kitabının tamamını formalize etti.
Automath'ta ortaya çıkan ``önermeleri tür olarak görme'' fikri, Howard'ın
çalışmalarıyla birlikte bu kitabın omurgasını oluşturan Curry--Howard
eşlemesinin kaynaklarından biridir.
\end{tarihce}

% ---------------------------------------------------------------------
\section{Lean 4'ü kurmak}\label{sec:kurulum}

\Lean{}'i denemek için hiçbir şey kurmanız gerekmez: Lean'in resmî sitesindeki
\emph{çevrimiçi editör}, tarayıcıda çalışır. Ancak kitaptaki projeyle
çalışmak için yerel bir kurulum öneriyoruz.

\begin{enumerate}[adimlar]
  \item \textbf{Visual Studio Code}'u kurun.
  \item VS Code'un eklenti mağazasından \textbf{Lean 4} eklentisini kurun.
  \item Eklentinin karşılama sayfasındaki adımları izleyin. Eklenti, Lean'in
        sürüm yöneticisi \lc{elan}'ı ve derleme aracı \lc{lake}'i sizin
        için kurar.
  \item Kitabın eşlik eden projesini indirip klasörünü VS Code'da açın ve
        bir terminalde şu komutu çalıştırın:
\begin{center}
\lc{lake build}
\end{center}
  \item \lc{ModernMantik/Bolum01.lean} dosyasını açın. Sağ tarafta
        \emph{Lean Infoview} paneli görünmelidir.
\end{enumerate}

\begin{dikkat}[Dikkat: Kurulum adımları değişebilir]
\Lean{} ve araçları hızla gelişiyor. Yukarıdaki adımlar bu kitabın yazıldığı
dönemde geçerlidir. Bir adım beklediğiniz gibi çalışmazsa, Lean'in resmî
sitesindeki güncel kurulum sayfasına başvurun. Eşlik eden projenin
\lc{lean-toolchain} dosyası, kitabın test edildiği Lean sürümünü belirtir.
\end{dikkat}

% ---------------------------------------------------------------------
\section{İlk adımlar: \#check ve \#eval}\label{sec:check-eval}

\Lean{} hem bir programlama dili hem de bir ispat dilidir. Bu ikili yapıyı
en iyi iki komutla görebiliriz: \lc{\textcolor{leancmd}{\bfseries\#eval}}
bir ifadeyi \emph{hesaplar}, \lc{\textcolor{leancmd}{\bfseries\#check}} ise
bir ifadenin \emph{türünü} söyler.

\leandosya[İlk komutlar]{Bolum01}{check_eval}

Son satır önemli bir şeyi gösteriyor: \lc{2 + 2 = 5} \emph{yanlış} bir
önermedir, ama \Lean{} onu itirazsız kabul eder ve türünün \leantype{Prop}
olduğunu söyler. Çünkü \leantype{Prop}, \emph{önermelerin} türüdür, yani
doğru ya da yanlış olabilen ifadelerin. Bir ifadenin önerme olması, doğru
olduğu anlamına gelmez.

\begin{onemlifikir}[Önemli Fikir: Bir önerme ile onun kanıtı farklı şeylerdir]
\lc{2 + 2 = 4} ve \lc{2 + 2 = 5} ifadelerinin ikisi de birer önermedir;
ikisinin de türü \leantype{Prop}'tur. Aralarındaki fark şudur: Birincisinin
bir \emph{kanıtını} inşa edebiliriz, ikincisinin edemeyiz. \Lean{}'de bir
önermenin doğru olduğunu söylemek, o önermenin bir kanıtını sunmak demektir.
\end{onemlifikir}

\Cref{tab:tur-benzetmesi}, bu fikrin \Lean{}'de nasıl yazıldığını
gösteriyor. Sol sütundaki tanıdık ifadelerle sağ sütundakiler arasında
yapısal bir benzerlik var.

\begin{table}[htbp]
\centering
\renewcommand{\arraystretch}{1.4}
\rowcolors{2}{tablosatir}{white}
\begin{tabular}{>{\raggedright\arraybackslash}p{4.2cm}
                >{\raggedright\arraybackslash}p{4.2cm}}
\rowcolor{tabloust}
\tablobaslik{Veri dünyası} & \tablobaslik{Mantık dünyası} \\
\lc{Nat} bir türdür                & \lc{P} bir önermedir (\lc{P : Prop}) \\
\lc{3 : Nat}                      & \lc{hp : P} \\
``3, bir doğal sayıdır''           & ``\lc{hp}, $P$'nin bir kanıtıdır'' \\
\lc{Nat} boş değildir              & $P$ doğrudur \\
\end{tabular}
\caption{Türler ve önermeler arasındaki benzerlik. Bu tablo, Curry--Howard
eşlemesinin ilk görünümüdür.}
\label{tab:tur-benzetmesi}
\end{table}

Benzetmenin son satırı en derin olanıdır. Bir türün bir elemanını
gösterebiliyorsak o tür ``dolu''dur; bir önermenin bir kanıtını
gösterebiliyorsak o önerme doğrudur. Bu fikri \cref{bol:onermeler1}'de
adım adım geliştirecek, \cref{bol:lambda}'te ise kesin bir teoreme
dönüştüreceğiz.

% ---------------------------------------------------------------------
\section{İlk teorem}\label{sec:ilk-teorem}

Şimdi ilk teoremimizi yazalım. Teoremimiz şunu söylüyor: \emph{Elimizde
$P$'nin bir kanıtı varsa, $P$ doğrudur.} Bu, olabilecek en basit teoremdir;
ama \Lean{}'in bir teoremi nasıl okuduğunu bütün parçalarıyla gösterir.

\leandosya[İlk teorem]{Bolum01}{ilk_teorem}

\Cref{fig:teorem-anatomisi}, bu tek satırın anatomisini gösteriyor.

\begin{figure}[htbp]
\centering
\begin{tikzpicture}[
    parca/.style={font=\leanfont\normalsize, inner sep=2pt},
    etiket/.style={font=\baslikfont\scriptsize, align=center, text width=2.6cm},
  ]
  \node[parca, text=leankw]   (a) at (0,0)    {\bfseries theorem};
  \node[parca]                (b) at (2.05,0) {ilk\_teorem};
  \node[parca]                (c) at (4.35,0) {(P : \textcolor{leantype}{\bfseries Prop})};
  \node[parca]                (d) at (6.85,0) {(\textcolor{leanhyp}{\bfseries hp} : P)};
  \node[parca]                (e) at (8.45,0) {: P};
  \node[parca]                (f) at (9.75,0) {:= \textcolor{leanhyp}{\bfseries hp}};

  \foreach \n/\r in {a/leankw, b/koyu, c/leantype, d/leanhyp, e/vurgu, f/ikincil}
    \draw[\r, thick, rounded corners=2pt] (\n.south west) rectangle (\n.north east);

  \node[etiket, text=leankw]   (ea) at (0,-1.35)   {anahtar kelime};
  \node[etiket, text=koyu]     (eb) at (2.05,1.35) {teoremin adı};
  \node[etiket, text=leantype] (ec) at (4.35,-1.35){``$P$ bir önerme olsun''};
  \node[etiket, text=leanhyp]  (ed) at (6.85,1.35) {varsayım: ``\lc{hp}, $P$'nin kanıtı''};
  \node[etiket, text=vurgu]    (ee) at (8.45,-1.35){iddia: ``$P$ doğrudur''};
  \node[etiket, text=ikincil]  (ef) at (10.2,1.35) {ispat: ``kanıt şudur: \lc{hp}''};

  \foreach \n/\m/\r in {a/ea/leankw, b/eb/koyu, c/ec/leantype,
                        d/ed/leanhyp, e/ee/vurgu, f/ef/ikincil}
    \draw[\r, -{Stealth[length=2mm]}] (\m) -- (\n);
\end{tikzpicture}
\caption{Bir \Lean{} teoreminin parçaları.}
\label{fig:teorem-anatomisi}
\end{figure}

Okunuşu şöyledir: ``$P$ bir önerme olsun ve elimizde $P$'nin \lc{hp} adlı
bir kanıtı bulunsun. O zaman $P$ doğrudur; kanıtı da \lc{hp}'nin
kendisidir.'' İki noktanın ardından gelen \lc{:=} işareti ``ispatı şudur''
diye okunur.

\subsection{Taktik modu}

Yukarıdaki ispat bir \emph{terimdir}: kanıtı doğrudan yazdık. Daha uzun
ispatlarda kanıtı tek seferde yazmak zordur. Bunun yerine \leankw{by}
anahtar kelimesiyle \emph{taktik moduna} geçer, ispatı adım adım inşa
ederiz:

\leandosya[Aynı teorem, taktik moduyla]{Bolum01}{ilk_teorem_taktik}

Taktik modunda \Lean{}, imlecin bulunduğu her noktada \emph{hedef durumunu}
Infoview panelinde gösterir (\cref{fig:infoview}). \leankw{by} kelimesinden
hemen sonra durum şöyledir:

\begin{hedefdurumu}[\texttt{by} sonrası]
P : Prop
hp : P
⊢ P
\end{hedefdurumu}

\noindent \lc{⊢} işaretinin üstündekiler elimizdekilerdir (\emph{bağlam}),
altındaki ise kanıtlamamız gereken şeydir (\emph{hedef}). \leantac{exact hp}
taktiği ``hedefin kanıtı tam olarak \lc{hp}'dir'' der ve ispatı bitirir:

\ispatbitti

\begin{figure}[htbp]
\centering
\begin{tikzpicture}[font=\leanfont\scriptsize]
  % Editör penceresi
  \draw[black!30, fill=white, rounded corners=4pt] (0,0) rectangle (7.6,4.2);
  \fill[black!8, rounded corners=4pt] (0,3.75) rectangle (7.6,4.2);
  \fill[black!8] (0,3.75) rectangle (7.6,3.95);
  \foreach \x/\c in {0.25/hataframe, 0.5/hedefframe, 0.75/bittiframe}
    \fill[fill=\c!80] (\x,3.975) circle (0.07);
  \node[anchor=west, font=\baslikfont\scriptsize, text=black!60] at (1.0,3.975) {Bolum01.lean};
  % Kod satırları
  \foreach \i/\t in {
     1/{\textcolor{leankw}{\bfseries theorem} ilk\_teorem' (P : \textcolor{leantype}{\bfseries Prop})},
     2/{\hspace{3.2em}(\textcolor{leanhyp}{\bfseries hp} : P) : P := \textcolor{leankw}{\bfseries by}},
     3/{\hspace{1em}\textcolor{leantac}{\bfseries exact} \textcolor{leanhyp}{\bfseries hp}}}
  {
    \node[anchor=west, text=black!35] at (0.15,{3.5-0.42*\i}) {\i};
    \node[anchor=west] at (0.55,{3.5-0.42*\i}) {\t};
  }
  % İmleç
  \fill[anarenk] (1.42,1.98) rectangle (1.46,2.28);
  \node[anchor=west, font=\baslikfont\tiny, text=anarenk] at (1.65,1.72) {imleç burada};

  % Infoview
  \draw[black!30, fill=hedefbg!60, rounded corners=4pt] (8.0,0) rectangle (13.4,4.2);
  \fill[hedefframe!30, rounded corners=4pt] (8.0,3.75) rectangle (13.4,4.2);
  \fill[hedefframe!30] (8.0,3.75) rectangle (13.4,3.95);
  \node[anchor=west, font=\baslikfont\scriptsize\bfseries, text=black!70] at (8.15,3.975) {Lean Infoview};
  \node[anchor=west, font=\baslikfont\scriptsize\bfseries] at (8.2,3.3) {{\sembolfont ▾} Tactic state};
  \node[anchor=west, text=black!60] at (8.3,2.85) {1 goal};
  \node[anchor=west] at (8.3,2.4) {P : \textcolor{leantype}{Prop}};
  \node[anchor=west] at (8.3,2.0) {\textcolor{leanhyp}{hp} : P};
  \node[anchor=west] at (8.3,1.55) {\textcolor{leansym}{⊢} P};
  \draw[hedefframe, thick, rounded corners=2pt] (8.2,1.3) rectangle (9.4,1.8);
  \node[anchor=west, font=\baslikfont\tiny, text=hedefframe!80!black] at (9.5,1.55) {hedef};
  \draw[black!40, decorate, decoration={brace, amplitude=4pt}]
       (10.3,2.6) -- (10.3,1.85) node[midway, right=5pt, font=\baslikfont\tiny, text=black!60] {bağlam};

  \draw[-{Stealth[length=2.5mm]}, thick, anarenk, dashed] (3.1,2.13) to[bend left=18] (8.15,2.6);
\end{tikzpicture}
\caption{VS Code'da editör ve Lean Infoview paneli. İmleç \leantac{exact}
satırına gelmeden önceki hedef durumu sağda görünür.}
\label{fig:infoview}
\end{figure}

\subsection{Hesaplamayla ispat}

Bazı önermeler doğrudan hesaplanarak doğrulanabilir. \lc{rfl} (yansıma,
İng.\ \emph{reflexivity}) ``iki taraf hesaplandığında aynıdır'' demektir:

\leandosya[Hesaplamayla ispat]{Bolum01}{rfl_ornek}

Aynı şeyi \lc{2 + 2 = 5} için denerseniz \Lean{} itiraz eder; çünkü iki taraf
hesaplandığında aynı değildir. Bir önermenin doğru olmasının ``\Lean{}'in
kanıtı kabul etmesi'' anlamına geldiğini burada açıkça görüyoruz.

% ---------------------------------------------------------------------
\section{Hata mesajlarını okumak ve \texttt{sorry}}\label{sec:hatalar}

İspat asistanıyla çalışmanın büyük kısmı, hata mesajlarını okumaktan oluşur.
Bunda utanılacak bir şey yoktur: Her hata mesajı, \Lean{}'in sizinle
konuştuğu bir andır. Bilerek yanlış bir ispat yazalım:

\begin{leankod}[Yanlış bir ispat (derlenmez)]
theorem yanlis (P Q : Prop) (hp : P) : Q := hp
\end{leankod}

\Lean{} \lc{hp}'nin altını kırmızıyla çizer ve aşağıdakine benzer bir mesaj
gösterir. (Mesajın tam ifadesi \Lean{} sürümüne göre biraz değişebilir.)

\begin{leanhata}[Lean hata mesajı]
type mismatch
  hp
has type
  P : Prop
but is expected to have type
  Q : Prop
\end{leanhata}

Okunuşu: ``\lc{hp} ifadesinin türü $P$, ama burada $Q$ türünden bir şey
beklenmektedir.'' Mantık diliyle: ``Sen bana $P$'nin kanıtını verdin, ama
ben $Q$'nun kanıtını istedim.'' \Lean{}, $P$'den $Q$'nun çıkacağına dair
hiçbir gerekçe göremediği için ispatı reddeder. Haklıdır da: $P$ doğru,
$Q$ yanlış olabilir.

\paragraph{sorry: ``Buraya sonra döneceğim''}
Bir ispatın bir kısmını henüz yazamıyorsanız, o kısmın yerine \leantac{sorry}
yazabilirsiniz. \Lean{} dosyanın geri kalanını denetlemeye devam eder,
ama \leantac{sorry} içeren her teoremi bir uyarıyla işaretler:

\begin{leankod}[Yarım kalmış bir ispat]
theorem henuz_bitmedi (P Q : Prop) (hp : P) : P ∧ Q := by
  sorry
\end{leankod}

\begin{dikkat}
\leantac{sorry}, bir ispatı \emph{kapatmaz}; yalnızca ertelemeyi sağlar.
\leantac{sorry} içeren bir teorem kanıtlanmış sayılmaz. Bir teoremin
gerçekten kanıtlanıp kanıtlanmadığını \lc{\textcolor{leancmd}{\bfseries\#print axioms}}
komutuyla denetleyebilirsiniz: \leantac{sorry} kullanılmışsa listede
\lc{sorryAx} görünür. Bu komuta \cref{bol:aksiyomlar}'da yeniden döneceğiz.
\end{dikkat}

Kitabın alıştırma dosyalarının hepsi bu yöntemle hazırlanmıştır: Teoremin
ifadesi verilir, ispatın yerinde \leantac{sorry} durur. Sizin göreviniz
\leantac{sorry}'yi gerçek bir ispatla değiştirmektir.

% ---------------------------------------------------------------------
\begin{arastirmanotu}
\textbf{Lean'in kısa geçmişi.} Lean, Leonardo de Moura tarafından Microsoft
Research'te başlatıldı; ilk sistem tanımı 2015'te yayımlandı
\citep{demoura2015lean}. \Lean{}, öncüllerinden farklı olarak büyük ölçüde
\emph{kendi dilinde} yazılmıştır; böylece kullanıcılar dili kendi taktikleri
ve notasyonlarıyla genişletebilir \citep{demoura2021lean4}. Topluluk
tarafından geliştirilen matematik kütüphanesi \Mathlib{}, bugün en büyük
birleşik formal matematik kütüphanelerinden biridir \citep{mathlib2020}.

\textbf{İleri okuma.} \Lean{}'e İngilizce giriş için resmî kaynaklar
\citet{tpil4} ve \citet{logicandproof}'tur. Bilgisayar bilimi
perspektifinden bir giriş için \citet{hitchhiker2026}'ya, programlama
yönü için \citet{fpil}'e bakılabilir. \Lean{}'in tip kuramının kesin
tanımı \citet{carneiro2019}'dadır.
\end{arastirmanotu}

\begin{bolumozeti}
\begin{itemize}
  \item Mantık, çıkarımların \emph{biçimine} bakarak geçerliliklerini
        inceler. Geçerlilik ile doğruluk farklı kavramlardır.
  \item İspat asistanları ispatları mekanik olarak denetler. \Lean{}'de
        güven, küçük bir çekirdeğe dayanır (de Bruijn ölçütü).
  \item \lc{\textcolor{leancmd}{\bfseries\#check}} türü,
        \lc{\textcolor{leancmd}{\bfseries\#eval}} değeri gösterir.
        \leantype{Prop}, önermelerin türüdür.
  \item \lc{hp : P}, ``\lc{hp}, $P$'nin bir kanıtıdır'' diye okunur.
        Bir önermeyi ispatlamak, onun türünden bir terim inşa etmektir.
  \item Taktik modunda (\leankw{by}) \Lean{}, her adımda bağlamı ve hedefi
        gösterir. \leantac{exact} hedefi kapatır, \leantac{sorry} erteler.
\end{itemize}
\end{bolumozeti}

\begin{alistirmalar}
\alistirma{1} Aşağıdaki çıkarımların her biri için geçerli olup
olmadığını ve sağlam olup olmadığını belirleyin:
(a) ``Bütün kediler memelidir. Tekir bir kedidir. O hâlde Tekir bir
memelidir.''
(b) ``Bütün balıklar yüzer. Yunus yüzer. O hâlde yunus bir balıktır.''

\alistirma{1} Aşağıdaki komutların çıktısını önce tahmin edin, sonra
\Lean{}'de deneyin:
\lc{\#check 3}, \lc{\#eval 3 * 7}, \lc{\#check (3 < 7)},
\lc{\#check "merhaba"}, \lc{\#eval "mer" ++ "haba"}.

\alistirma{1} \lc{ilk\_teorem}'in ispatındaki \lc{hp}'yi silip yerine
\leantac{sorry} yazın. Infoview'da ne değişiyor? Sonra \lc{\#print axioms}
komutunu çalıştırın.

\alistirma{2} Aşağıdaki teoremi ispatlayın. Elinizdeki varsayımlardan
hangisi gereksiz?
\begin{center}
\lc{theorem ikinci (P Q : Prop) (hp : P) (hq : Q) : Q := sorry}
\end{center}

\alistirma{2} \cref{sec:hatalar}'teki \lc{yanlis} teoremini \Lean{}'de
yazın. Hata mesajını kendi cümlelerinizle açıklayın. Sonra varsayımlara
tek bir varsayım ekleyerek teoremi ispatlanabilir hâle getirin.

\alistirma{3} (Tartışma) de Bruijn ölçütü neden önemlidir? Bir ispat
asistanının elaboratöründe bir hata bulunduğunu düşünün. Bu hata hangi
koşullarda yanlış bir teoremin kabul edilmesine yol açabilir, hangi
koşullarda açamaz?
\end{alistirmalar}
